\section*{Chapter \ref{ch:qm:optimal-stopping-and-impulse-control} --- Optimal Stopping and Impulse Control}

\begin{solution}{exo:qm:optimal-stopping-and-impulse-control:1}
$b^* = (6 \times 300 \times 400/0.5)^{1/4} = 34.6$ million; $C^* = 2\sqrt{300 \times 400 \times 0.5/6} = 200$ dollars a day; $\sigma^2/b^{*2} = 1/3$ hedge
a day, one every three days.
\end{solution}

\begin{solution}{exo:qm:optimal-stopping-and-impulse-control:2}
$U_2 = G_2 = (2, 0, 0)$ at $S_2 = 102, 100, 98$. $U_1 = \max(1, 1) = 1$ at 101 and $\max(0, 0) = 0$ at 99; $U_0 = \max(0, 0.5) = 0.5$. At time 0
continue ($U > G$); at time 1 stopping is optimal (at 101 exactly indifferent, at 99 nothing is lost).
\end{solution}

\begin{solution}{exo:qm:optimal-stopping-and-impulse-control:3}
No: whether it is the largest of the month is known only at the month's end, so it is not a \omterm{def:qm:probability-at-speed:stopping}{stopping time}, and no policy can hedge
there.
\end{solution}

\begin{solution}{exo:qm:optimal-stopping-and-impulse-control:4}
Doubling $K$: band $\times 2^{1/4} = 1.19$ (41.2 million), cost $\times\sqrt2$ (283 dollars). Doubling $\sigma$: band $\times\sqrt2$ (49.0 million), cost $\times 2$
(400 dollars).
\end{solution}

\begin{solution}{exo:qm:optimal-stopping-and-impulse-control:5}
$b^* = (3 \times 25 \times 400/2)^{1/3} = 24.7$ million; cost $25 \times 400/49.3 + 0.5 \times 24.7^2/3 = 304$ dollars a day.
\end{solution}

\begin{solution}{exo:qm:optimal-stopping-and-impulse-control:6}
Below the threshold $V(x) = Ae^{\theta x}$ (the solution of $\mathcal LV = rV$ that stays bounded as $x \to -\infty$); value matching $Ae^{\theta b} = b - c$ and
\omterm{def:qm:optimal-stopping-and-impulse-control:fbp}{smooth pasting} $A\theta e^{\theta b} = 1$ give $\theta(b - c) = 1$. With the numbers, $\theta = (-0.5 + \sqrt{0.25 + 0.16})/4 = 0.0351$ and $b^* = 5 + 28.5 =
33.5$ bp.
\end{solution}

\begin{solution}{exo:qm:optimal-stopping-and-impulse-control:7}
197.5 dollars a day at 390 checks and 199.6 at 1\,560, against 200. Checked discretely, the exposure overshoots the band before a
hedge, so cycles last a little longer and there are fewer tickets; the extra risk charge of the overshoot is smaller than the fees
saved, and the gap closes like $\sqrt{\Delta t}$.
\end{solution}

\begin{solution}{exo:qm:optimal-stopping-and-impulse-control:8}
At $b = 10$ the policy costs $300 \times 400/100 + 0.5 \times 100/6 = 1\,208$ dollars a day, six times the optimum: the fees of hedging a dozen
times a day dwarf the risk they remove. Risk and fees must be traded off; less risk is not free.
\end{solution}

\begin{solution}{pb:qm:optimal-stopping-and-impulse-control:1}
\textbf{1.} $C(b) = K\sigma^2/b^2 + \gamma b^2/6$.
\textbf{2.} $b^* = (6K\sigma^2/\gamma)^{1/4} = 34.6$ million.
\textbf{3.} 200 dollars a day.
\textbf{4.} Every three days on average.
\textbf{5.} $300 + 0.5 \times 400/2 = 400$ dollars a day: the band saves 50\%.
\textbf{6.} 29.1 million (16\% narrower) and 141 dollars a day.
\textbf{7.} 49.0 million and 400 dollars a day.
\textbf{8.} Half: at the optimum fees and risk charge are equal.
\textbf{9.} 197.5 and 199.6 dollars a day.
\textbf{10.} Discrete checks let the exposure overshoot the band, lengthening cycles and saving tickets.
\textbf{11.} 24.7 million and 304 dollars a day.
\textbf{12.} Trigger 42.4 million, reset 12.8 million, 418 dollars a day.
\textbf{13.} Every million traded costs 25 dollars, so stopping short of zero saves proportional cost, and the next cycle, starting
closer to the band, is still long enough to amortise the fixed fee.
\textbf{14.} \omterm{def:qm:optimal-stopping-and-impulse-control:singular}{Singular control} trades infinitesimally, only when the exposure touches the band, with total traded amount finite; \omterm{def:qm:optimal-stopping-and-impulse-control:impulse}{impulse
control} makes finitely many discrete trades, each paying the fixed fee.
\textbf{15.} $b^* = 33.5$ bp in closed form; the daily grid stops at 32.2 bp and values the position at zero at 8.79 bp against 8.80.
\textbf{16.} As a variance charge, yes: the P\&L variance of an open exposure is proportional to $X^2$. A charge proportional to $|X|$ (a
VaR-type limit) changes the exponents of the band laws.
\textbf{17.} Wider: flows that offset themselves make waiting cheaper, since part of the exposure disappears without a trade.
\textbf{18.} Correlated exposures partly offset; banding the net risk hedges less often than banding each currency.
\textbf{19.} Named result: \emph{the hedging band}: the optimal no-trade band is $\pm(6K\sigma^2/\gamma)^{1/4} = \pm 34.6$ million, costing 200 dollars a day,
half the cost of daily hedging; the band grows like the fourth root of the ticket fee (16\% narrower if the fee halves).
\textbf{20.} Inside the band the risk removed by a hedge is worth less than the fee it costs, and waiting keeps the option to hedge a
larger exposure for the same fee.
\end{solution}

\begin{solution}{iq:qm:optimal-stopping-and-impulse-control:1}
A band: do nothing while the exposure is within $\pm b$, trade back to zero at the edge, with $b = (6K\sigma^2/\gamma)^{1/4}$ for a fee $K$, flow
volatility $\sigma$ and risk charge $\gamma X^2$. It equalises fees and risk charge.
\iqlookfor{\omterm{def:qm:optimal-stopping-and-impulse-control:notrade}{no-trade region}, the trade-off, and the scaling.}
\end{solution}

\begin{solution}{iq:qm:optimal-stopping-and-impulse-control:2}
At the optimal boundary the value of waiting and of stopping agree (value matching) and so do their slopes. If the value of waiting
met the payoff at an angle, moving the boundary slightly would raise the value: the threshold would not be optimal.
\iqlookfor{first-order optimality of the boundary.}
\end{solution}

\begin{solution}{iq:qm:optimal-stopping-and-impulse-control:3}
With a fixed fee, the cost per unit time is fee times hedge frequency $\sigma^2/b^2$ plus a risk charge $\propto b^2$: $b^4 \propto K$. With a
proportional cost, trading happens at rate $\sigma^2/b$ in size at the boundary, so the cost is $c\sigma^2/b$ plus $\propto b^2$: $b^3 \propto c$.
\iqlookfor{how often and how much one trades as a function of $b$.}
\end{solution}

\begin{solution}{iq:qm:optimal-stopping-and-impulse-control:4}
The smallest \omterm{def:qm:probability-at-speed:martingale}{supermartingale} dominating the reward process, computed backward as the maximum of stopping now and the expected
value of continuing; the first time it equals the reward is optimal. The price of an American option is the \omterm{def:qm:optimal-stopping-and-impulse-control:snell}{Snell envelope} of its
discounted payoff under the \omterm{def:qm:girsanov-and-changes-of-numeraire:emm}{risk-neutral measure} (One Quant Book~5, chapter~6).
\iqlookfor{the backward recursion and its optimality.}
\end{solution}

\begin{solution}{iq:qm:optimal-stopping-and-impulse-control:5}
Trade to a reset level strictly inside the band, never to the band itself, so a hedge cannot immediately trigger another;
measure the exposure net of pending hedges; debounce noisy exposure feeds; cap the number of hedges per interval; log each
decision with the exposure and the band in force; compute the band from current inputs, not a constant.
\iqlookfor{hysteresis, in-flight orders, and auditability.}
\end{solution}

\begin{solution}{iq:qm:optimal-stopping-and-impulse-control:6}
$f(x) = \E_x\int_0^\tau X^2\,dt$ solves $\tfrac12\sigma^2f^{\prime\prime} = -x^2$ with $f(\pm b) = 0$, so $f(x) = (b^4 - x^4)/(6\sigma^2)$ and $f(0) = b^4/(6\sigma^2)$.
\iqlookfor{Dynkin's formula and a boundary-value problem.}
\end{solution}
